\begin{titlepage}
\centering

{\Large\bfseries DSLO v0.7 — Geometry Examples \& Demonstrations\par}
\vspace{1.5em}

{\large\bfseries Geometric expressions in drift, collapse, recovery, and lawful manifold transitions.\par}
\vspace{2em}

{\Large\bfseries DSLO v0.7 --- Geometry\par}
\vspace{2em}

{\large Donald Slowicki\par}
{\small Inda Moment Incorporated\par}

\vspace{3em}

{\small
Document Version: v0.7\\
Substrate Version: DSLO v0.7\\
Discipline Version: DISCO v0.7\\
Release Date: August 2026\\
Static Release ID: SM-DSLO-07027\\
}

\vspace{3em}

{\small
© 2026 Inda Moment Incorporated • 
DSLO™ and Signal Ecology™ are trademarks of Inda Moment Incorporated • 
\href{https://www.tnopsi.com}{www.tnopsi.com} • 
\href{https://sites.google.com/tnopsi.com/dslo-protocol/defensive-publication}{DSLO Defensive Publication} • 
DOI: \href{https://zenodo.org/records/21864440}{10.5281/zenodo.21864440}
}

\end{titlepage}


\begin{abstract}
DSLO geometry spans six classes expressed across biological, cultural, physical, abstract, and 
machine-facing systems. This paper provides one expanded example for each class—Morphological Mimicry 
\& Adaptive Geometry, Distributed Cognition \& Non-Neural Intelligence, Cultural Signal Systems, 
Natural Optimization Phenomena, Abstract \& Universal Geometry Instances, and Modern \& Machine-Facing 
Geometry Instances. Each example includes a class overview, a representative instance, and an outline 
of how DSLO geometric primitives---drift, load, pressure, collapse thresholds, recovery attractors, and 
lawful manifold transitions---appear within that system.

The examples demonstrate how local geometric behavior reflects broader DSLO substrate invariants, 
showing how drift accumulates, how load and pressure shape manifold transitions, how collapse emerges 
under structural or operational limits, and how recovery attractors restore stability across diverse 
substrates. Together, these examples provide a compact cross-surface demonstration of DSLO geometry, 
illustrating invariant patterns that hold across biological mimicry, distributed non-neural intelligence, 
cross-cultural signal encoding, natural optimization processes, abstract mathematical partitioning, and 
modern engineered systems.

Full definitions, invariants, schemas, and classification structures for all geometry classes and 
instances are available in the DSLO v0.7 repository, which provides the canonical substrate-level 
specification for the DSLO geometry framework.
\end{abstract}


\section{Introduction}
\label{sec:introduction}

The DSLO v0.7 substrate defines a unified geometric framework for describing drift, load, pressure,
collapse thresholds, recovery attractors, and lawful manifold transitions across a wide range of
systems. These geometric primitives appear consistently in biological, cultural, ecological, abstract,
and machine-facing environments, forming a universal structure that is independent of substrate and
implementation. DSLO geometry provides a common language for understanding how systems move across
manifold states, how collapse emerges under structural or operational limits, and how recovery
attractors restore stability.

This paper presents a compact set of examples illustrating how DSLO geometry manifests across six major
classes: Morphological Mimicry \& Adaptive Geometry, Distributed Cognition \& Non-Neural Intelligence,
Cultural Signal Systems, Natural Optimization Phenomena, Abstract \& Universal Geometry Instances, and
Modern \& Machine-Facing Geometry Instances. Each class is represented by one expanded instance:
Boquila leaf mimicry, slime-mold network optimization, Incan Quipu encoding, river-delta branching,
Voronoi partition geometry, and internet routing geometry. These examples demonstrate how DSLO
geometric primitives appear in real and synthetic systems, revealing invariant patterns that hold
across biological mimicry, distributed non-neural intelligence, cross-cultural signal encoding, natural
optimization processes, abstract mathematical structures, and engineered computational environments.

Each section provides a class overview, a representative instance, and an outline of how drift, load,
pressure, collapse, and recovery appear within that system. The examples show how local geometric
behavior reflects broader DSLO substrate invariants, and how lawful transitions guide systems from
unstable to stable manifold states. A summary section consolidates the geometry classes, the example
structure, and the universal primitives shared across all six surfaces.

These examples are not exhaustive; rather, they serve as accessible entry points into the broader
classification structures defined in the DSLO v0.7 repository. Full definitions, invariants, schemas,
and registry entries are available in the main DSLO v0.7 release, which provides the canonical
substrate-level specification for all geometry classes and their associated instances.


\section{Morphological Mimicry \& Adaptive Geometry}
\label{sec:morphological-mimicry}

This section presents a representative instance from the class of Morphological Mimicry \& Adaptive 
Geometry. The example illustrates how local morphological changes express DSLO geometric primitives 
such as drift, load, pressure, collapse thresholds, recovery attractors, and lawful manifold transitions.

\subsection{Class Overview}
Morphological Mimicry \& Adaptive Geometry refers to systems that alter their visible form, structure, 
or presentation in response to environmental signals. These systems exhibit geometry that shifts under 
drift and load, often producing adaptive configurations that reduce collapse pressure. The class includes 
biological mimicry, adaptive camouflage, and signal-responsive morphological changes.

\subsection{Representative Instance}
A representative instance for this class is \emph{Boquila trifoliolata}, a plant capable of mimicking the 
leaf shape of nearby host species. Boquila adjusts its morphology based on environmental cues, producing 
leaf geometries that closely resemble surrounding foliage. This instance demonstrates adaptive geometry 
driven by signal intake and environmental drift.

\subsection{Geometric Expression}
In Boquila, drift appears as gradual morphological deviation toward the geometry of nearby leaves. Load 
accumulates as environmental pressure to reduce predation or increase camouflage effectiveness. Pressure 
emerges when the plant must match complex or rapidly changing leaf geometries. Collapse corresponds to 
failure to mimic effectively, resulting in increased visibility or ecological disadvantage. Recovery occurs 
when the plant re-stabilizes its morphology toward a successful mimicry attractor, restoring geometric 
alignment with its environment.

\subsection{Class--Instance Relationship}
Boquila expresses the broader class geometry by demonstrating how morphological systems respond to 
environmental drift and load through adaptive shape transitions. The instance shows how local geometry 
can shift across manifold states while maintaining lawful transitions that preserve ecological function. 
Its behavior exemplifies the class-wide pattern of signal-driven morphological adaptation.

\subsection{Reference to Full Registry}
For full definitions, invariants, and classification structures associated with Morphological Mimicry \& 
Adaptive Geometry, refer to the DSLO v0.7 registry and geometry classification surfaces in the main 
DSLO v0.7 release.

\section{Distributed Cognition \& Non-Neural Intelligence}
\label{sec:distributed-cognition}

This section provides an example from Distributed Cognition \& Non-Neural Intelligence, illustrating how 
DSLO geometry appears in systems that compute, coordinate, or optimize without centralized control. 
The example demonstrates how distributed agents express drift, load, pressure, collapse thresholds, 
recovery attractors, and lawful manifold transitions through collective behavior.

\subsection{Class Overview}
Distributed Cognition \& Non-Neural Intelligence refers to systems in which problem-solving, optimization, 
or decision-making emerges from the interaction of many simple components rather than from a single 
centralized controller. These systems exhibit geometric behavior distributed across agents, pathways, 
or nodes, with drift and load accumulating through local interactions. Collapse and recovery occur 
through global reconfiguration of the distributed manifold.

\subsection{Representative Instance}
A representative instance for this class is \emph{Physarum polycephalum} (slime mold), a non-neural 
organism capable of solving shortest-path problems, optimizing networks, and adapting its structure 
to environmental constraints. Slime mold forms dynamic transport networks that reorganize in response 
to drift, load, and pressure, demonstrating distributed geometric intelligence.

\subsection{Geometric Expression}
In slime mold, drift appears as gradual reallocation of mass toward nutrient sources or efficient 
pathways. Load accumulates when multiple competing nutrient sites or obstacles increase the complexity 
of the transport network. Pressure emerges when the organism must maintain connectivity under resource 
constraints or environmental perturbations. Collapse corresponds to the failure of inefficient or 
overloaded pathways, which are pruned or abandoned. Recovery occurs when new, more efficient pathways 
form, stabilizing the network into a lower-cost attractor geometry.

\subsection{Class--Instance Relationship}
Slime mold expresses the broader class geometry by demonstrating how distributed systems reorganize 
their structure through local interactions that collectively produce global optimization. The instance 
shows how drift, load, pressure, collapse, and recovery propagate across a distributed manifold, 
revealing lawful transitions between network configurations. Its behavior exemplifies the class-wide 
pattern of decentralized geometric intelligence.

\subsection{Reference to Full Registry}
For full definitions, invariants, and classification structures associated with Distributed Cognition 
\& Non-Neural Intelligence, refer to the DSLO v0.7 registry and geometry classification surfaces in 
the main DSLO v0.7 release.

\section{Cultural Signal Systems (Cross-Cultural)}
\label{sec:cultural-signal-systems}

This section introduces a cross-cultural signal system example, demonstrating how DSLO geometry appears 
in human-designed encoding systems that transmit structured information across time, space, and social 
contexts. The example illustrates cultural drift, collapse, recovery, and lawful transitions within 
signal geometry.

\subsection{Class Overview}
Cultural Signal Systems refer to human or cross-cultural structures that encode, store, and transmit 
information through geometric or symbolic forms. These systems exhibit drift as signals change across 
generations or contexts, load as complexity or usage increases, pressure as interpretive or structural 
constraints accumulate, collapse when meaning or structure fails, and recovery when stable attractors 
re-emerge through cultural reinforcement or re-standardization.

\subsection{Representative Instance}
A representative instance for this class is the \emph{Incan Quipu}, a knot-based encoding system used 
for record-keeping, communication, and administrative coordination. Quipus encode numerical and 
categorical information through knot geometry, fiber arrangement, and positional structure. This 
instance demonstrates how cultural systems use geometric encoding to maintain meaning across distributed 
contexts.

\subsection{Geometric Expression}
In the Quipu system, drift appears as gradual variation in knot patterns, fiber colors, or positional 
conventions across regions or generations. Load accumulates as the system must encode increasingly 
complex administrative or cultural information. Pressure emerges when interpretive ambiguity increases 
or when structural conventions diverge. Collapse corresponds to loss of interpretability or breakdown 
of shared encoding standards. Recovery occurs when communities re-align knot geometry, re-establish 
interpretive norms, or reinforce stable encoding attractors through cultural practice.

\subsection{Class--Instance Relationship}
The Quipu expresses the broader class geometry by demonstrating how cultural systems maintain meaning 
through structured geometric encoding. The instance shows how drift, load, pressure, collapse, and 
recovery operate within cultural manifolds, revealing lawful transitions between stable and unstable 
encoding states. Its behavior exemplifies the class-wide pattern of cross-cultural signal geometry.

\subsection{Reference to Full Registry}
For full definitions, invariants, and classification structures associated with Cultural Signal Systems, 
refer to the DSLO v0.7 registry and geometry classification surfaces in the main DSLO v0.7 release.

\section{Natural Optimization Phenomena}
\label{sec:natural-optimization}

This section presents a natural optimization example, showing how DSLO geometric primitives manifest in 
physical or ecological systems. The example illustrates how drift, load, pressure, collapse thresholds, 
and recovery attractors appear in environments governed by physical constraints and energy minimization.

\subsection{Class Overview}
Natural Optimization Phenomena refers to systems in which geometric structure emerges from physical 
processes that minimize energy, reduce transport cost, or optimize flow. These systems exhibit drift as 
conditions change, load as resources accumulate or redistribute, pressure as constraints intensify, 
collapse when structural limits are exceeded, and recovery when the system reorganizes into a stable 
attractor geometry. Examples include branching networks, fluid flows, crystal formation, and other 
self-organizing physical structures.

\subsection{Representative Instance}
A representative instance for this class is \emph{river-delta branching}, a geometry formed by the 
interaction of sediment flow, water velocity, and environmental constraints. River deltas develop 
branching networks that distribute flow across multiple channels, optimizing transport while responding 
to drift in sediment deposition and load from upstream sources.

\subsection{Geometric Expression}
In river-delta branching, drift appears as gradual shifts in sediment distribution and channel geometry. 
Load accumulates as increased water volume or sediment concentration pushes the system toward new 
branching configurations. Pressure emerges when existing channels approach their transport limits or 
when environmental constraints restrict flow. Collapse corresponds to channel failure, blockage, or 
redirection, often producing abrupt geometric transitions. Recovery occurs when new channels form or 
existing ones reorganize, stabilizing the delta into a lower-energy attractor geometry that restores 
flow distribution.

\subsection{Class--Instance Relationship}
River-delta branching expresses the broader class geometry by demonstrating how physical systems 
optimize structure under changing conditions. The instance shows how drift, load, pressure, collapse, 
and recovery propagate through a physical manifold, revealing lawful transitions between branching 
states. Its behavior exemplifies the class-wide pattern of natural optimization driven by physical 
constraints and energy minimization.

\subsection{Reference to Full Registry}
For full definitions, invariants, and classification structures associated with Natural Optimization 
Phenomena, refer to the DSLO v0.7 registry and geometry classification surfaces in the main DSLO v0.7 
release.

\section{Abstract \& Universal Geometry Instances}
\label{sec:abstract-universal}

This section provides an abstract or universal geometry instance, illustrating invariant geometric 
behavior independent of substrate. The example demonstrates how DSLO geometric primitives appear in 
mathematical or structural systems that do not rely on biological, cultural, or physical embodiment.

\subsection{Class Overview}
Abstract \& Universal Geometry Instances refer to geometric structures that arise from formal rules, 
mathematical constraints, or universal partitioning behavior. These systems exhibit drift as parameters 
change, load as structural constraints accumulate, pressure as boundaries tighten, collapse when 
geometric stability fails, and recovery when the system reorganizes into a stable attractor. Because 
these geometries are substrate-independent, they reveal DSLO primitives in their purest form.

\subsection{Representative Instance}
A representative instance for this class is the \emph{Voronoi pattern}, a universal geometric structure 
formed by partitioning space into regions based on proximity to a set of points. Voronoi diagrams appear 
in mathematics, physics, biology, and engineered systems, making them an ideal example of substrate-neutral 
geometry.

\subsection{Geometric Expression}
In Voronoi geometry, drift appears as gradual movement of generating points, causing region boundaries 
to shift. Load accumulates when additional points are introduced or when spatial constraints tighten. 
Pressure emerges when regions become highly irregular or when point density increases. Collapse occurs 
when boundaries destabilize or when partitioning fails to maintain coherent regions. Recovery happens 
when the system re-stabilizes into a new Voronoi attractor, restoring lawful partitioning under updated 
conditions.

\subsection{Class--Instance Relationship}
Voronoi patterns express the broader class geometry by demonstrating how universal structures respond to 
changes in parameters through lawful transitions. The instance shows how drift, load, pressure, collapse, 
and recovery operate in purely abstract manifolds, revealing DSLO geometry independent of substrate. 
Its behavior exemplifies the class-wide pattern of invariant geometric response under rule-based 
constraints.

\subsection{Reference to Full Registry}
For full definitions, invariants, and classification structures associated with Abstract \& Universal 
Geometry Instances, refer to the DSLO v0.7 registry and geometry classification surfaces in the main 
DSLO v0.7 release.


\section{Modern \& Machine-Facing Geometry Instances}
\label{sec:machine-facing}

This section introduces a modern or machine-facing geometry example, demonstrating how DSLO geometric 
primitives appear in engineered, computational, or networked systems. The example illustrates drift, 
load, pressure, collapse, recovery, and lawful manifold transitions in environments shaped by machine 
constraints and operational dynamics.

\subsection{Class Overview}
Modern \& Machine-Facing Geometry Instances refer to systems whose geometric behavior emerges from 
engineered architectures, computational processes, or machine-level operations. These systems exhibit 
drift as operational conditions shift, load as demand or traffic increases, pressure as constraints 
tighten, collapse when system limits are exceeded, and recovery when the system reorganizes into a 
stable attractor state. Examples include routing networks, distributed compute systems, and engineered 
control surfaces.

\subsection{Representative Instance}
A representative instance for this class is \emph{internet routing geometry}, the dynamic structure 
formed by routing protocols, network topology, and traffic flow. Routing geometry evolves as packets 
move through distributed nodes, responding to drift in network conditions and load from traffic 
patterns. This instance demonstrates how engineered systems express DSLO geometry through operational 
reconfiguration.

\subsection{Geometric Expression}
In internet routing geometry, drift appears as gradual changes in latency, link quality, or node 
availability. Load accumulates as traffic increases or as routing tables expand. Pressure emerges when 
critical links approach saturation or when routing decisions become unstable due to rapid changes. 
Collapse corresponds to routing failure, congestion collapse, or partitioning of the network. Recovery 
occurs when routing protocols re-stabilize paths, redistribute load, or converge on new attractor 
geometries that restore connectivity and reduce operational stress.

\subsection{Class--Instance Relationship}
Internet routing geometry expresses the broader class behavior by demonstrating how machine-facing 
systems reorganize under drift, load, and pressure through lawful transitions between operational 
states. The instance shows how collapse and recovery propagate across engineered manifolds, revealing 
DSLO geometric primitives in computational and networked environments. Its behavior exemplifies the 
class-wide pattern of machine-level geometric adaptation.

\subsection{Reference to Full Registry}
For full definitions, invariants, and classification structures associated with Modern \& Machine-Facing 
Geometry Instances, refer to the DSLO v0.7 registry and geometry classification surfaces in the main 
DSLO v0.7 release.


\section{Summary of Geometry Classes and Example Structure}
\label{sec:summary}

\subsection{Overview}
This paper presents one representative instance for each of the six DSLO geometry classes. 
Each example demonstrates how DSLO geometric primitives---drift, load, pressure, collapse thresholds, 
recovery attractors, and lawful manifold transitions---appear across biological, cultural, physical, 
abstract, and machine-facing systems. The examples serve as compact demonstrations of the broader 
classification structures defined in the DSLO v0.7 repository.

\subsection{Geometry Classes}
The six geometry classes covered in this paper are:

\begin{enumerate}
    \item \textbf{Morphological Mimicry \& Adaptive Geometry} \\
    Systems that alter form or presentation in response to environmental signals.

    \item \textbf{Distributed Cognition \& Non-Neural Intelligence} \\
    Systems that exhibit collective or decentralized problem-solving behavior.

    \item \textbf{Cultural Signal Systems (Cross-Cultural)} \\
    Human or cultural systems that encode, transmit, and transform structured signals.

    \item \textbf{Natural Optimization Phenomena} \\
    Physical or ecological systems that exhibit optimization geometry under drift and load.

    \item \textbf{Abstract \& Universal Geometry Instances} \\
    Geometry that appears independent of substrate, often in mathematical or structural forms.

    \item \textbf{Modern \& Machine-Facing Geometry Instances} \\
    Engineered or computational systems that express DSLO geometry under operational constraints.
\end{enumerate}

\subsection{Outline Structure}
Each section follows a consistent outline:

\begin{itemize}
    \item \textbf{Class Overview} --- A brief description of the geometry class.
    \item \textbf{Representative Instance} --- A selected example expressing the class geometry.
    \item \textbf{Geometric Expression} --- Notes on drift, collapse, recovery, and manifold transitions.
    \item \textbf{Class--Instance Relationship} --- How the instance reflects the broader class geometry.
    \item \textbf{Reference to Full Registry} --- Pointer to the DSLO v0.7 classification structures.
\end{itemize}

\subsection{Universal Geometry}
Across all six classes, DSLO geometry exhibits a universal structure:

\begin{itemize}
    \item \textbf{Drift} --- Gradual deviation under signal or load.
    \item \textbf{Load} --- Accumulated pressure on the system.
    \item \textbf{Pressure} --- Threshold approach toward collapse.
    \item \textbf{Collapse} --- Structural or functional failure region.
    \item \textbf{Recovery} --- Return to stability via attractor geometry.
    \item \textbf{Lawful Transitions} --- Constrained movement between manifold states.
\end{itemize}

These primitives appear consistently across biological, cultural, physical,
abstract, and machine-facing systems, demonstrating the universality of DSLO
geometry. Full definitions, invariants, and classification structures are
available in the DSLO v0.7 repository.

\medskip

\href{https://github.com/Signal-Ecology/DSLO-v0.7-Semantic-Substrate-Specification/tree/main/EXAMPLES}
     {DSLO v0.7 Example Geometry Library}




\end{document}
